\documentclass[10pt]{article}
\usepackage[T1]{fontenc}
\usepackage{lmodern}
\usepackage[margin=0.72in]{geometry}
\usepackage{amsmath,amssymb}
\usepackage{microtype}
\usepackage[hidelinks]{hyperref}
\usepackage{enumitem}

\begin{document}
\begin{center}
{\Large\bfseries The Restricted Invertibility Conjecture in\\
arXiv:1010.6090 Is Already Answered\par}
\smallskip
{\small Literature-status note --- August 17, 2026}
\end{center}
\small

\paragraph{Status.}
\texttt{literature\_already\_answered}.  This is an exact status
identification, not a new mathematical result.

\paragraph{Source problem.}
Nikolski and Vasyunin state the following conjecture in Section~3.2,
PDF page~19 \cite{NV}:
\begin{quote}
For every bounded operator $T$ on a Hilbert space $H$ and every orthogonal
basis $(e_j)$ satisfying
$\inf_j\|Te_j\|/\|e_j\|>0$, there is a finite partition
$\bigcup_{s=1}^r I_s=\mathbb N$ such that every restriction
$T|_{H_{I_s}}$ is left invertible.
\end{quote}
Immediately before the display, the source records the known implication
\[
 \text{positive Kadison--Singer solution}\quad\Longrightarrow\quad
 \text{RIC},
\]
citing Casazza--Christensen--Lindner--Vershynin.

\paragraph{Answer.}
Marcus, Spielman, and Srivastava prove two theorems known to imply a positive
solution of the Kadison--Singer problem \cite[abstract and Corollary~1.5]{MSS}.
Their introduction explicitly includes the Bourgain--Tzafriri conjecture
among the equivalent formulations of Kadison--Singer.  Combining their
positive solution with the implication already recorded by the source gives
a positive answer to the displayed RIC.

The supporting authors were aware of this relationship: arXiv:1306.3969
names the Bourgain--Tzafriri conjecture in its equivalence list rather than
merely supplying an unrelated theorem later recognized to apply.

\paragraph{Other source signal.}
Remark~3 of arXiv:1010.6090 says that the first Blaschke product constructed
there is not known to have a noninvertible element exactly at its threshold.
Theorem~2.5 later in the same paper constructs a refined product with exactly
that endpoint property.  This signal is a same-paper ask-and-answer passage,
not a separate open problem.

\paragraph{Scope limitation.}
The paper also names stronger variants for Schauder bases (B-RIC) and
summation bases (SB-RIC), and refutes SB-RIC.  This note classifies only the
formally displayed RIC.  It makes no claim about the current status of B-RIC.

\begin{thebibliography}{9}
\bibitem{NV}
N. Nikolski and V. Vasyunin,
\emph{Invertibility threshold for $H^\infty$ trace algebras, and effective
matrix inversions}, arXiv:1010.6090 (2010), especially PDF p.~19.

\bibitem{MSS}
A. W. Marcus, D. A. Spielman, and N. Srivastava,
\emph{Interlacing Families II: Mixed Characteristic Polynomials and the
Kadison--Singer Problem}, arXiv:1306.3969v4; Ann. of Math. 182 (2015),
327--350.
\end{thebibliography}

\end{document}
